\paragraph{The Hebrew Bible and Early Christianity.}
The stories of Joseph (Gen 40–41) and Daniel (Dan 2) established the dream as a legitimate medium of revelation.  
In Late Antiquity and the medieval Christian world, patristic writers debated whether dreams were divine or demonic, while popular “dream manuals” circulated widely despite ecclesiastical suspicion \citep{miller1994lateantiquity}.

\paragraph{Classical Islam}
Hadith traditions classify dreams into three types—true dreams (\textit{ru’yā}), confused ones (\textit{adghāth ahlām}), and satanic ones.  
Vast oneirocritical handbooks were compiled, the most famous attributed to Ibn Sīrīn (7th–8th century, though authorship is complex).  
Philosophers such as al-Kindī and al-Fārābī also wrote naturalistic explanations of dreaming \citep{sirr2000islam,ibnsirin2007}.

\paragraph{Medieval Christian Europe and Byzantium}
Collections like the \textit{Somniale Danielis} spread alphabetized glossaries of dream symbols with “good/bad” moral valuations and tests of authenticity; manuscripts survive in both Latin and vernaculars \citep{cappozzo2018somniale}.